\newpage
\chapter{CONCLUSION \& FUTURE WORK}
\pagestyle{fancy}

% ================================================
% 6.1 Conclusions
% ================================================
\section{Conclusions}

This thesis presented \textbf{CA-DQStream}, a context-aware framework for streaming data quality monitoring built on Apache Kafka and Apache Flink. The framework addresses three critical limitations of existing batch-oriented and rule-based systems: inability to adapt to distributional shifts over time, inability to detect multivariate anomalies requiring inter-feature relationships, and inability to provide context-sensitive validation criteria across heterogeneous data segments.

CA-DQStream was validated on 72M NYC Yellow Taxi records (14 months real TLC data + 12 months simulated data). Four core contributions were evaluated:

\begin{enumerate}
    \item \textbf{Online Autoencoder + Memory Module:} Denoising autoencoder ($34$D $\rightarrow$ $76$D $\rightarrow$ $34$D, trained with noise $\sigma=0.001$) combined with a FIFO memory buffer ($N=8{,}192$ vectors, production configuration with $4\times$ margin above the ablation-identified optimal of $N=2{,}048$). On the 1M-record evaluation set, CA-DQStream achieves AUC-PR 0.903, AUC-ROC 0.997, F1 0.860, and FPR 0.65\% (below the 5\% production target), substantially outperforming Isolation Forest (F1 0.432, FPR 2.89\%) and the MemStream baseline (F1 0.825, FPR 0.73\%).

    \item \textbf{4D Context-Aware Thresholding:} Threshold matrix ($80$ cells in production: $10$ neighborhood zones $\times$ $8$ temporal context cells) with automated 95th-percentile computation. The 4D formulation (trip\_type $\times$ time\_window $\times$ day\_type $\times$ neighborhood, $588$ cells) achieves marginally higher performance on the full-data evaluation but suffers from sparse cells with fewer than 100 training samples. The production 2D instantiation captures the two most impactful dimensions (spatial and temporal), consistent with empirical variance reduction ranking. ContextBeta thresholds provide a marginal but consistent improvement over global thresholds (+0.001 AUC-PR).

    \item \textbf{Staged Validation Pipeline with Early Exit:} Four-layer pipeline (Schema Validation $\rightarrow$ Canary Rules $\rightarrow$ Online Autoencoder $\rightarrow$ MetaAggregator) with conditional early exit routing. Records failing lightweight stages bypass ML inference. Context-aware routing reduces RocksDB I/O by approximately 90\% via Bloom filter pre-check. Per-record scoring throughput on the production Kafka-Flink cluster: $1{,}687{,}817$ records evaluated in $643$ seconds with memory footprint $\approx 56$ MB per TaskManager slot.

    \item \textbf{Intelligent Evolution Controller:} ADWIN-based drift detection on quality meta-streams. \textbf{Limitation:} Drift detection claims (detection latency, retrain frequency, adaptation strategy distribution) are validated exclusively on months 1--14 of real NYC TLC data (Jan 2024 -- Feb 2025). Months 15--26 are GMM-simulated without novel drift patterns and are excluded from drift evaluation.
\end{enumerate}

All four research questions are supported by experimental evidence: CA-DQStream significantly outperforms Isolation Forest and the streaming MemStream baseline on AUC-PR and FPR at $p < 0.001$ (Wilcoxon signed-rank test, Bonferroni-corrected). The grid search confirms memory size $N$ as the most impactful hyperparameter; AUC-PR improves from 0.339 ($N=16$) to 0.916 ($N=2{,}048$) and degrades for $N > 4{,}096$. Feature engineering is the single largest component contributor (+0.87 AUC-PR over raw features). The context-aware threshold matrix provides a modest but consistent improvement (+0.001 AUC-PR).

% ================================================
% 6.2 Limitations
% ================================================
\section{Limitations}

Three key limitations constrain the current framework:

\begin{enumerate}
    \item \textbf{Single-domain evaluation:} CA-DQStream is validated exclusively on the NYC Yellow Taxi dataset. Generalizability to other streaming domains (financial transactions, IoT sensors, healthcare data) requires domain-specific parameter tuning and context dimension redesign.

    \item \textbf{Synthetic anomaly ground truth:} Production anomalies are unlabeled. Evaluation relies on synthetic injection, which may not capture the full distribution of real-world anomalies (e.g., coordinated fraud patterns spanning multiple trips).

    \item \textbf{No anomaly explanation:} The autoencoder produces anomaly scores but does not inherently explain \emph{why} a record is flagged. Feature attribution and counterfactual analysis remain future work.
\end{enumerate}

% ================================================
% 6.3 Future Work
% ================================================
\section{Future Work}

Four directions extend from this work:

\begin{enumerate}
    \item \textbf{Multi-domain validation:} Extend evaluation to financial transactions (credit card fraud with merchant-category and geography context), IoT sensor streams (manufacturing quality control), and healthcare data (ICU vital signs with patient-specific baselines).

    \item \textbf{Anomaly explanation:} Augment anomaly scores with reconstruction error attribution to identify which input features drive the score. Implement per-context explanation dashboards to support operator root-cause analysis.

    \item \textbf{Approximate nearest-neighbor indexing:} Replace brute-force k-NN over the 8,192-vector memory buffer with HNSW or FAISS to support larger memory capacities and higher throughput (e.g., $>$100K events/sec).

    \item \textbf{Deep learning variants:} Explore Variational Autoencoders for probabilistic anomaly scoring, LSTM for temporal anomaly patterns, and attention-weighted memory access for improved pattern matching.
\end{enumerate}

% ================================================
% 6.4 Broader Impact
% ================================================
\section{Broader Impact}

Beyond taxi trip data quality, CA-DQStream contributes three transferable insights to data engineering and MLOps:

\begin{itemize}
    \item \textbf{Zero-interruption model updates:} Broadcast State enables production ML systems to update models without pausing stream processing, applicable to real-time recommendation, fraud detection, and predictive maintenance pipelines.
    \item \textbf{Learned vs. manual DQ thresholds:} Demonstrating that DQ rules can be learned from data rather than manually specified, reducing DevOps burden while providing quantitative evidence for DQ decisions.
    \item \textbf{Regulatory compliance:} Audit trails (drift events logged to MinIO), explainable business rules, and Grafana dashboard exports support compliance reporting for finance, healthcare, and government domains.
\end{itemize}

\vspace{1.5cm}
\begin{flushright}
\textit{Le Dac Thinh} \\
\textit{Ha Noi, May 2025}
\end{flushright}
